% Transparent, Post-Quantum Anonymous Set Membership from FRI-STARKs over BabyBear and Mersenne-31
% LaTeX export of the ePrint draft (auto-converted from the Markdown source; repo metadata stripped).
% Compiles with pdflatex (article class + standard packages). To target IACR formatting, swap
% \documentclass{article} for {iacrcc} or {llncs} and adjust the title/author block accordingly.
\documentclass[11pt]{article}

\usepackage{lmodern}        % scalable Latin Modern fonts (needed for microtype expansion)
\usepackage[T1]{fontenc}
\usepackage[utf8]{inputenc}
\usepackage[margin=1in]{geometry}
\usepackage{amsmath,amssymb,amsthm}
\usepackage{booktabs}
\usepackage{enumitem}
\usepackage{microtype}
\usepackage{xurl}           % allow long URLs to break anywhere (fixes overfull lines in the bibliography)
\usepackage[hidelinks]{hyperref}

\setlength{\emergencystretch}{3em}  % soak up tiny (<2pt) overfull lines
\setlength{\tabcolsep}{4pt}         % tighter table columns so the wide tables fit

\newcommand{\GF}[1]{\mathrm{GF}(#1)}
\newcommand{\cat}{\mathbin{\|}}          % field-sequence concatenation separator
\newcommand{\wout}{w_{\mathrm{out}}}

\title{Transparent, Post-Quantum Anonymous Set Membership\\
from FRI-STARKs over BabyBear and Mersenne-31}
\author{Enoch Kuskoff\\ Enkom Tech\\ \texttt{en.cipher@enkom.tech}}
\date{}

\begin{document}
\maketitle

\begin{center}\itshape
A construction-and-analysis note that solicits external review of its security arguments.
\end{center}

\begin{quote}
\textbf{Scope of this note.} A self-contained description of a transparent anonymous set-membership
argument and its security analysis, presented as a generic primitive. The building blocks
(FRI-STARKs, Poseidon, sponge hashing, nullifier-based membership) are standard; the contribution is
the \emph{concrete instantiation and its parameters} together with an explicit security analysis. We
do \textbf{not} claim the construction is proven secure: two central questions, the
$\GF{p^2}$-state Poseidon round counts (O1, specific to Arm A) and the formal zero-knowledge
simulator (O4, shared by both arms), are stated as \textbf{open problems on which we solicit review}
(Sections 6--7), not as established results (O2 and O3 reduce to the permutation's security, i.e.\ O1).
\end{quote}

\begin{abstract}
We describe and analyze a transparent, plausibly post-quantum argument of \emph{anonymous set
membership with a nullifier}: a prover demonstrates, in zero knowledge, that it knows a secret $t$
whose leaf $L = H(t)$ lies in a public Merkle accumulator with root $\mathit{root}$, and
simultaneously publishes a context-bound nullifier $N = H(\mathit{domain} \cat t \cat \mathit{ctx})$
that is unlinkable across contexts but stable within one. The argument is a FRI-STARK over a small
$31$-bit prime field, using a single algebraic hash (a Poseidon-family permutation in a wide-sponge
mode) for the leaf, every Merkle node, and the nullifier. Transparency (no trusted setup) and a
hash-only soundness assumption make the construction conservative against quantum adversaries.

We give a complete security analysis along four axes: (O1) algebraic / Gr\"obner-basis security of
the permutation round counts; (O2) collision and second-preimage resistance of the truncated-output
sponge; (O3) domain separation between the leaf and nullifier pre-images; and (O4) the zero-knowledge
simulator for the hiding-PCS instantiation. We \textbf{implement, measure, and compare} two field
instantiations of the hash. Arm A runs the permutation \emph{state} over the quadratic extension
$\GF{p^2}$ of Mersenne-31; Arm B runs over the base field of BabyBear with the deployed Poseidon2.
Both are implemented, and their proof-system soundness is computed to $128$-bit post-quantum
(Section 6.5); at equal security, \textbf{Arm B's proofs are $\approx 14$--$18\times$ smaller and
verify $\approx 16$--$17\times$ faster}, and we recommend it. We explicitly solicit review of the two
residuals our engineering arguments do not settle: the extension-field round-count parameterization
\emph{specific to Arm A} (O1), and the zero-knowledge simulator's masking-degree budget (O4, shared
by both arms).

\smallskip
\noindent\textbf{Keywords:} anonymous set membership $\cdot$ nullifiers $\cdot$ zero-knowledge proofs
$\cdot$ FRI-STARKs $\cdot$ transparent setup $\cdot$ post-quantum cryptography $\cdot$
Poseidon\,/\,Poseidon2 $\cdot$ algebraic hashing $\cdot$ BabyBear\,/\,Mersenne-31.
\end{abstract}

\section{Introduction}

Anonymous set membership with nullifiers underlies anonymous signaling, one-person-one-vote /
rate-limiting, private credentials, and privacy-preserving messaging (e.g.\ Semaphore-style signaling
\cite{Sem}, rate-limiting nullifiers \cite{RLN}, and privacy pools). The standard shape is: members
are accumulated in a Merkle tree over a commitment $H(t)$ of a secret $t$; to act, a member proves in
zero knowledge that its leaf is in the tree and reveals a deterministic nullifier $N$ derived from
$t$ and a context, so that double-action \emph{within} a context is detectable while actions
\emph{across} contexts are unlinkable.

Most deployed instantiations use pairing- or discrete-log-based SNARKs (Groth16, PLONK) with a
trusted or universal setup, and rely on assumptions broken by a quantum adversary. We target a
\textbf{transparent} (setup-free) and \textbf{plausibly post-quantum} instantiation by realizing the
relation as an AIR proven with a \textbf{FRI-STARK} \cite{BBHR18,BGKS20}, whose soundness rests only
on collision resistance of a hash and the FRI low-degree test. The single algebraic hash is a
Poseidon-family permutation \cite{GKRRS21,GKS23} in a wide-sponge mode, chosen for low AIR degree.

A second post-quantum route is \emph{lattice-based} succinct ZK, e.g.\ LaBRADOR \cite{LaBRADOR} and
the recent toolkit of Biasioli et al.\ \cite{BBL26} that adds zero-knowledge to it, which yields
markedly smaller proofs (under $100$\,KB) but rests on structured-lattice assumptions
(Module-SIS\,/\,LWE). We deliberately take the more conservative \emph{hash-only} path (its sole
assumption is collision resistance and it needs no trusted setup), and target a purpose-built
membership relation rather than a general prover (Section 5).

\paragraph{Contributions.}
\begin{enumerate}[leftmargin=*]
\item A self-contained description of the membership AIR: trace layout, in-circuit sponge constraints,
  Merkle path threading, the nullifier sub-circuit, and the public-input binding (Sections 3--4).
\item Two concrete, \textbf{implemented and measured} field instantiations of the hash: a
  $\GF{p^2}$-\emph{state} variant over Mersenne-31 (Arm A) and a base-field Poseidon2 variant over
  BabyBear (Arm B), with their digest widths, security margins, wire sizes, and a side-by-side cost
  comparison (Section 5).
\item A four-axis security analysis (O1--O4, Section 6) carried out against the published round-count
  and sponge-indifferentiability bounds, including an adversarial self-check against the recent
  algebraic-attack literature.
\item A precise statement of the two residual questions on which we seek external review (Section 7).
\end{enumerate}

We make no claim that the underlying techniques are new. The value we hope to add is a \emph{published,
reviewable} concrete parameter set and analysis for a transparent, hash-only anonymous-membership
argument over a small prime field.

\section{Preliminaries}

\paragraph{Field.} Arm A uses $p = 2^{31} - 1$ (Mersenne-31), a 31-bit prime; $\GF{p^2}$ denotes the
quadratic extension realized as \texttt{Complex<GF(p)>} with $i^2 = -1$ (valid since $-1$ is a
non-residue: $p \equiv 3 \pmod 4$); a $\GF{p^2}$ element packs two base limbs, $\log_2|\GF{p^2}|
\approx 62$. Arm B (Section 5.2) instead uses the \textbf{BabyBear} prime $q = 2^{31} - 2^{27} + 1$
(also $31$-bit, but two-adic of order $2^{27}$) as its base field. Field-generic statements below are
written over $F$.

\paragraph{Algebraic permutation.} We use a Poseidon-family permutation $\pi$ over $F^t$ (state width
$t$) with an $x^\alpha$ S-box: $\alpha = 5$ for Arm A ($\gcd(5, p-1) = \gcd(5, p^2-1) = 1$, so $x^5$
permutes both $\GF{p}$ and $\GF{p^2}$); $\alpha = 7$ for Arm B (since $5 \mid q-1$ makes $x^5$
non-bijective over BabyBear). There are $R_F$ full rounds and $R_P$ partial rounds, ARC
(add-round-constant), and an MDS / Poseidon2 external$+$internal linear layer. Round constants are
derived by an XOF / Grain-LFSR from a fixed domain string.

\paragraph{Sponge.} $H$ is the sponge \cite{BDPV07} over $\pi$ with rate $r$, capacity $c$
($t = r + c$), \texttt{10*1} padding, and a \emph{truncated} output equal to the first $\wout$ cells
of the final state. We use $H$ for the leaf $L = H(t)$, every Merkle node
$H(\mathit{left} \cat \mathit{right})$, and the nullifier $N = H(\mathit{domain} \cat t \cat
\mathit{ctx})$.

\paragraph{FRI-STARK / AIR.} The relation is expressed as an Algebraic Intermediate Representation
(AIR): an execution trace (a matrix over the field) and a set of polynomial constraints of bounded
degree, proven low-degree after the relation is reduced to a quotient via FRI \cite{BBHR18}. We use a
STARK whose trace, constraints, and FRI all live over a common field $F$ (Section 5 discusses
$F = \GF{p^2}$ for Arm A vs.\ the two-adic base-field BabyBear configuration for Arm B; both run on a
radix-2 FRI domain, not a circle-STARK \cite{HLP24}).

\paragraph{Zero-knowledge.} Plain STARK traces are not zero-knowledge. We obtain ZK via a \emph{hiding}
polynomial commitment: the trace is low-degree-extended at an increased blow-up with
vanishing-polynomial randomization, Merkle commitments are salted with fresh CSPRNG output, and FRI
runs on the randomized codeword \cite{Hab22,BCRSVW19}. Only the public statement
$(\mathit{root}, \mathit{ctx}, N)$ is revealed.

\section{The Membership Relation}

Fix a domain constant $\mathit{domain}$ (a baked circuit constant, \emph{not} a witness or public
input), derived as the first $|\mathit{domain}|$ cells of $H(S)$ for a fixed statement string $S$.
The argument proves, for public $(\mathit{root}, \mathit{ctx}, N)$:
\[
\begin{aligned}
\exists\,(t, \mathit{path}) :\quad
  & \textsf{MerkleVerify}(\mathit{root},\, L = H(t),\, \mathit{path}) = \textsf{true} \\
  & \wedge\ N = H(\mathit{domain} \cat t \cat \mathit{ctx}) \\
\text{reveal only } & (\mathit{root}, \mathit{ctx}, N); \quad L \text{ and } t \text{ remain secret.}
\end{aligned}
\]
$\mathit{root}$ is the verifier-supplied, trusted accumulator root; $\mathit{ctx}$ is a public
per-context label; $N$ is the nullifier. Membership soundness reduces to: a prover cannot make the
AIR accept unless it folds a genuine $H(t)$ to the trusted $\mathit{root}$ and emits the
correctly-derived $N$. Anonymity reduces to: the proof, $\mathit{root}$, and $N$ reveal nothing about
which leaf $L$ was used (within the anonymity set of the tree) and $N$ is unlinkable across distinct
$\mathit{ctx}$.

\section{Construction (the AIR)}

\paragraph{Trace.} One row per Merkle level, plus row-0 carrying the leaf and nullifier sponge blocks.
The public values are the $\wout$-cell $\mathit{root}$, the $|\mathit{ctx}|$-cell $\mathit{ctx}$, and
the $\wout$-cell $N$.

\paragraph{In-circuit sponge.} Each hash invocation is unrolled into $\lfloor L/r \rfloor + 1$
permutation applications ($L = $ input length in cells). Every round intermediate is pinned by
exactly one constraint: ARC, the $x^\alpha$ S-box, and the linear layer in full rounds; in partial rounds
the non-active lanes are copy-constrained so they cannot be freely chosen. Capacity is threaded as
the linear-layer \emph{expression} between permutations (no free capacity column), and the digest is
the first $\wout$ cells of the genuine final-state expression, never a free read-back column. The
absorb $+$ \texttt{10*1} padding schedule matches a value-level reference, checked across input
lengths.

\paragraph{Merkle threading.} Per level: a boolean $\mathit{dir}$ selects the node input order
($(\mathit{running}, \mathit{sibling})$ vs.\ $(\mathit{sibling}, \mathit{running})$);
$\mathit{running}$ is updated to the node hash; a transition constraint forces
$\mathit{running}' = \mathit{parent}$ (next-row $\mathit{running}$); the last row binds
$\mathit{parent} = \mathit{root}$. Trees of varying depth are handled by deterministically extending
$\mathit{root}$ with zero-sibling levels (padding-blind in the AIR; soundness rests on the verifier
feeding its \emph{own} trusted $\mathit{root}$).

\paragraph{Nullifier.} The $|\mathit{domain}|$-cell $\mathit{domain}$ prefix is injected as constant expressions (no
committed domain column). The pre-image is exactly $\mathit{domain} \cat t \cat \mathit{ctx}$. The
single committed $t$ slice feeds \emph{both} the leaf sponge and the nullifier pre-image (so leaf and
nullifier are structurally same-$t$). Row-0 gated constraints bind $\mathit{running} = L$,
$\mathit{ctx}$ to the public $\mathit{ctx}$, and the nullifier output to the public $N$.

\paragraph{Degree.} Max constraint degree is $\alpha$ (the ungated $x^\alpha$ S-box: $5$ for Arm A, $7$ for Arm B); selectors gate only
degree-1 bindings. The quotient-chunk count is derived from the constraint AST by both prover and
verifier, so any degree overflow fails at prove time.

A structural audit of this AIR (sponge under-constraint, Merkle threading, leaf / same-$t$, nullifier
integrity, padding rows, padding-root, public binding, degree / FRI config) found no relation a
malicious prover can exploit to accept a proof that violates the decoded statement, \emph{modulo} the
cryptographic assumptions O1--O4 below.

\section{Field Instantiation}

The single design lever with security and cost consequences is the \textbf{field over which the hash
permutation state runs}. We describe, implement, and measure both arms. \textbf{Both are realized,
and their proof-system soundness computes to $128$-bit post-quantum (Section 6.5); Arm B is the
instantiation we recommend.} Arm A runs the permutation state over the quadratic extension
$\GF{p^2}$ of Mersenne-31; Arm B runs over the base field of BabyBear with the deployed Poseidon2.
Measured at equal security, Arm B's proofs are $\approx 14$--$18\times$ smaller and verify
$\approx 16$--$17\times$ faster (Section 5.4), and its round-count obligation (O1) is a near-citation
of a standard parameter set rather than the novel extension-field analysis Arm A requires. All other
aspects of the construction (Sections 3--4) are identical between the two.

\subsection{Arm A: \texorpdfstring{$\GF{p^2}$}{GF(p2)}-state (the analyzed instantiation)}

State width $t = 7$, rate $r = 2$, capacity $c = 5$, $R_F = 8$, $R_P = 60$, $\alpha = 5$, over
$\GF{p^2}$. Each state cell is a $\GF{p^2}$ element ($\approx 62$ bits). Digest $\wout = 5$ cells
$= 5 \times 8 = 40$ bytes $= 5 \times 62 \approx 310$ bits. The STARK trace, constraints, and FRI all
run over $\GF{p^2}$.
\begin{itemize}[leftmargin=*]
\item \textbf{Pro:} the wide $\approx 62$-bit cells give a $5$-cell, $310$-bit digest and $155$-bit
  capacity collision resistance with only $c = 5$.
\item \textbf{Con:} running the \emph{whole} trace and FRI over $\GF{p^2}$ pays a $\approx 2\times$
  field tax on all prover arithmetic and committed data; and the published round-count formulas
  (Section 6, O1) target the \emph{base} prime field, so the extension-field state requires the
  dedicated analysis of O1.
\end{itemize}
The FRI challenges are drawn from a \textbf{degree-3 extension $\GF{p^6} \approx 186$ bits}: the
$\approx 62$-bit value field is too small to serve as its own challenge field at $128$-bit soundness,
so the membership prover / verifier lift the challenge / OOD sampling into $\GF{p^6}$ (Section 6.5).

Public statement: $\mathit{root}(40) \cat \mathit{ctx}(16) \cat N(40) = 96$ bytes; per-cell encoding
$4$-byte real $\cat$ $4$-byte imaginary, canonical.

\subsection{Arm B: base-field BabyBear / Poseidon2 (the realized, recommended instantiation)}

State width $t = 16$, $R_F = 8$ ($4 + 4$), $R_P = 13$, $\alpha = 7$, over the \textbf{BabyBear} prime
field $q = 2^{31} - 2^{27} + 1$ ($\log_2 q \approx 30.9$), using the \emph{deployed} Poseidon2
parameters \cite{GKS23,Plonky3}: the external block-circulant $4 \times 4$ MDS layer, the optimized
internal diagonal, and Grain-LFSR round constants from the reference generator, i.e.\ the width-16
instance shipped by Plonky3 / SP1. BabyBear, not Mersenne-31, is the base field for a deliberate
reason: M31's base field has 2-adicity $1$ (no radix-2 FFT subgroup), so a base-M31 STARK would
require \emph{circle} STARKs \cite{HLP24}; BabyBear's 2-adicity $27$ gives a native two-adic FFT
domain, so Arm B runs on the \textbf{same \texttt{TwoAdicFriPcs}} machinery as Arm A; only the field
and permutation change, not the proof system. BabyBear's S-box is $x^7$ ($\alpha = 7$, since
$5 \mid q-1$ makes $x^5$ non-bijective).

To preserve a $\geq 256$-bit digest with $\approx 31$-bit cells, the capacity and output widen to
$\lceil 256/30.9 \rceil = 9$ cells: sponge $t = 16,\ r = 7,\ c = 9,\ \wout = 9$ ($\approx 278$-bit
capacity $\to \approx 139$-bit collision; $9 \times 4 = 36$-byte digest). FRI challenges are drawn
from a degree-5 extension $\GF{q^5} \approx 155$ bits, sized for $128$-bit soundness (Section 6.5).
\begin{itemize}[leftmargin=*]
\item \textbf{Pro:} O1 reduces from a novel extension-field analysis to a near-\emph{citation} of the
  standard, widely-deployed BabyBear Poseidon2 parameter set (Plonky3 / SP1); and the base-field
  trace ($4$-byte cells) avoids the $\GF{p^2}$ field tax across the whole proof. Measured, the effect
  is decisive (Section 5.4).
\item \textbf{Realization:} Arm B is fully implemented and tested \textbf{for functional correctness}:
  the field, the deployed Poseidon2 (KAT against the published reference vector), the in-circuit
  gadget, the wide sponge / Merkle, the membership AIR, and a transparent \textbf{and} zero-knowledge
  prover / verifier, with an exhaustive per-column under-constraint audit and a negative-proof matrix
  over all sub-AIRs. This is functional evidence, \textbf{not} a soundness proof: like Arm A, Arm B
  remains open on O1 / O4 (Sections 6--7). Public statement:
  $\mathit{root}(36) \cat \mathit{ctx}(16) \cat N(36) = 88$ bytes, per-cell $4$-byte canonical
  little-endian.
\end{itemize}

\subsection{What the field choice does and does not change}

The field choice governs O1 (round-count security) and, through it, O2 / O3 (collision and
domain-separation, which assume a secure permutation). It does \textbf{not} affect O4 (the
zero-knowledge simulator), which is a property of the hiding PCS and the masking-degree budget, not
of the hash field. A reviewer's verdict on O4 transfers unchanged between the two arms.

\subsection{Measured cost (both arms, equal \texorpdfstring{$128$}{128}-bit-PQ security)}

Both arms are implemented and benchmarked at their $128$-bit-PQ configurations (Arm A: $\GF{p^6}$
challenge, \texttt{log\_blowup}~3; Arm B: $\GF{q^5}$ challenge, \texttt{log\_blowup}~4; both bind on a
SHAKE256 Merkle commitment at $128$ bits; Section 6.5). Single-thread, median of 5, membership depth
32:

\begin{table}[h]
\centering\footnotesize
\begin{tabular}{@{}lrrr@{}}
\toprule
 & Arm A ($\GF{p^2}$ / Poseidon256) & Arm B (BabyBear / Poseidon2) & Arm B advantage \\
\midrule
Trace row width          & $17\,152$         & $1\,661$         & $\approx 10\times$ narrower \\
Bytes per trace row       & $137\,216$        & $6\,644$         & $\approx 21\times$ smaller  \\
Proof size, transparent   & $\approx 16.4$ MB & $\approx 1.04$ MB & $\approx 16.6\times$ smaller \\
Proof size, zero-knowledge & $\approx 16.7$ MB & $\approx 1.21$ MB & $\approx 14.4\times$ smaller \\
Verify time               & $\approx 510$ ms  & $\approx 30$ ms   & $\approx 17\times$ faster    \\
\bottomrule
\end{tabular}
\end{table}

The gap is \textbf{at equal security} (both reach the same $128$-bit hash floor), not a
size-for-security trade. It tracks bytes-per-trace-row: FRI opens $\approx (\#\mathrm{queries}) \times
(\text{trace width})$ field elements, and Arm A's Poseidon256 ($t = 7$, $68$ rounds, $8$-byte
$\GF{p^2}$ cells) yields a far wider, heavier row than Arm B's Poseidon2 ($t = 16$, $21$ rounds,
$4$-byte cells). This field / permutation effect, not the FRI blowup, dominates, so it holds across
depths and in both transparent and zero-knowledge modes.

\section{Security Analysis (O1--O4)}

The structural audit shows the constraints faithfully encode the intended hash and Merkle relations.
The following are the cryptographic obligations that an automated / structural argument cannot
discharge. We give engineering derivations against the published formulas; we do \textbf{not} claim
these are a substitute for independent expert review (Section 7).

\subsection*{O1: Round-count / algebraic-degree security}

\textbf{Instance (Arm A).} $\alpha = 5$, $t = 7$, $R_F = 8$, $R_P = 60$ ($68$ rounds) over $\GF{p^2}$,
$\log_2|F| \approx 62$. Target $M = 128$-bit.

We evaluate the published Poseidon round-count bounds \cite{GKRRS21} (\S5.5), \cite{ABM23}:
statistical, interpolation, and the Gr\"obner-basis family, reproduced from the canonical
generators. With $n = \log_2 p$ the field-size term, $\log_5 2 \approx 0.43068$:
\begin{itemize}[leftmargin=*]
\item \textbf{Statistical:} $R_F \geq 6$ if $M \leq \lfloor n - (\alpha-1)/2 \rfloor (t+1)$, else
  $R_F \geq 10$.
\item \textbf{Interpolation:} $R_F + R_P \geq 1 + \lceil \log_5 2 \cdot \min(M,n) \rceil +
  \lceil \log_5 t \rceil$.
\item \textbf{Gr\"obner:} several bounds, incl.\ $(t-1)(R_F + R_P) \geq (t-2) + M/(2 \log_2 \alpha)$
  and the Macaulay-style $\mathrm{cost}_{\mathrm{gb}} \geq M$.
\item \textbf{Margin:} $+2$ full rounds and $+7.5\%$ partial rounds over the minimum.
\end{itemize}

\begin{table}[h]
\centering\footnotesize
\begin{tabular}{@{}llll@{}}
\toprule
Regime & Binding bound (interpolation) & Margined optimal $(R_F, R_P)$ & Deployed $8 + 60$ \\
\midrule
$n = 62$ (realistic $\GF{p^2}$)   & $1 + \lceil 0.43068 \cdot 62 \rceil + 2 = 30$ & $(8, 26)$, total 34 & $\approx 2\times$ total / $2.3\times R_P$ \\
$n = 31$ (pessimistic subfield)   & $1 + \lceil 0.43068 \cdot 31 \rceil + 2 = 17$ & $(8, 14)$, total 22 & $\approx 4.3\times R_P$ \\
\bottomrule
\end{tabular}
\end{table}

$\mathrm{cost}_{\mathrm{gb}}$ for $(8, 60)$ evaluates to $\approx 349$ bits $\gg 128$. So $8 + 60$
clears every published bound under both the realistic and worst-case-subfield substitutions.

\textbf{Extension-field subtlety (the crux for review).} $\GF{p^2}$'s tower structure could
\emph{reduce} security only via a \textbf{subfield descent}: if $\GF{p}^{\,7}$ were invariant under
the round function, an attacker could work in the 31-bit subfield ($\min(M,n) \to 31$) and unlock
invariant-subspace / subspace-trail Gr\"obner attacks \cite{Bariant22,GKR25}. The defense is
\textbf{generic round constants}: drawing both the real and imaginary limb of every constant
independently from the XOF gives each constant a pseudorandom nonzero imaginary part
($\Pr[\mathrm{imag} = 0] \approx 2^{-31}$ per constant), so a single off-subfield constant per round
maps $\GF{p}^{\,7}$ off itself and breaks subfield invariance. \emph{Caveat for the reviewer:} if the
linear (MDS) layer is built from base-field entries, the linear layer alone preserves $\GF{p}^{\,7}$;
the subfield defense then rests \textbf{entirely} on the constants' genericity. A future change that
de-randomized the constants' imaginary parts would re-expose the subfield.

\textbf{Adversarial self-check.} The 2024--2026 algebraic-attack wave (the FreeLunch / CheapLunch
Gr\"obner attacks \cite{FreeLunch,CheapLunch} and subspace-trail Gr\"obner cryptanalysis of Poseidon
\cite{GKR25}) pressures $R_P$ by linearizing partial rounds; $R_P = 60$ is $2.3$--$4.3\times$ the
requirement, and the instances those attacks render borderline run $R_P$ at $\approx 1\times$. The
off-envelope concern reduces to the constant-genericity check. Notably, \cite{GKR25} finds that the
\emph{original} Poseidon round-count model can itself mis-estimate the rounds required for some
instantiations (under- or over-counting); this makes the large $R_P = 60$ margin reassuring but is a
further reason expert review, not formula evaluation alone, is warranted. The case does not survive
automated scrutiny, but the residual against \emph{future} cryptanalysis of a $\GF{p^2}$-\emph{state}
permutation (an unusual deployment choice) is real, and is why we seek review.

\textbf{Arm B contrast.} Arm B's permutation is the deployed BabyBear Poseidon2 \cite{GKS23,Plonky3}:
$\alpha = 7$, $t = 16$, $R_F = 8$, $R_P = 13$ over the base prime field ($\log_2 q \approx 30.9$),
the \textbf{standard, widely-deployed} instantiation (Plonky3 / SP1) the round-count formulas target
directly. O1 becomes confirmation that these published parameters are used unchanged (a base-field
citation with no tower / subfield-invariance / off-envelope-state concern) rather than the novel
extension-field analysis Arm A requires.

\subsection*{O2: Collision / second-preimage of the truncated sponge}

For a sponge on a random permutation, collision and (second-)preimage resistance are, up to the
capacity \cite{BDPV08},
\[
\min\!\big(2^{c \cdot \log_2|F| / 2},\ 2^{\wout \cdot \log_2|F| / 2}\big).
\]
\begin{itemize}[leftmargin=*]
\item \textbf{Arm A:} capacity $5 \cdot 62 = 310$ bits $\to 2^{155}$; output $5 \cdot 62 = 310$ bits
  $\to 2^{155}$ and $\geq 256$-bit digest. Required capacity for $M = 128$ is
  $c \geq 2M / \log_2|F| = 256/62 \approx 4.13$; $c = 5$ clears it. Truncation to
  $\wout = 5\ (\geq 4.13)$ preserves $\geq 128$-bit output collision resistance; \emph{internal}
  collisions (which enable Merkle-node equivocation) are governed by the capacity, not the output
  width, so truncation does not weaken node-collision resistance. Sponges with nonzero capacity are
  not length-extendable.
\item \textbf{Arm B:} capacity $9 \cdot 30.9 \approx 278$ bits $\to \approx 2^{139}$; $9$-cell
  $\approx 278$-bit digest $\geq 256$ bits (BabyBear $\approx 31$-bit cells; $c = 9 \geq
  2M / \log_2 q \approx 8.3$ clears the $128$-bit floor).
\end{itemize}
O2 holds \emph{given} O1 (the permutation is a secure PRP at the $128$-bit level).

\subsection*{O3: Domain separation (no leaf\texorpdfstring{$\leftrightarrow$}{<->}nullifier cross-collision)}

The leaf preimage is the secret $t$; the nullifier preimage is the strictly longer
$\mathit{domain} \cat t \cat \mathit{ctx}$; both use the same sponge with \texttt{10*1} padding.
(i)~\emph{Cross-protocol:} a distinct $\mathit{domain}$ constant forces a distinct absorption from
round 1, giving disjoint nullifier images (standard prefix domain separation). (ii)~\emph{No
leaf$\leftrightarrow$nullifier impersonation:} the two preimages differ in length (so they absorb in
a different number of permutations), and the nullifier carries the $\mathit{domain}$ prefix the leaf
lacks; a collision between them is therefore a generic cross-length collision (O2), not a structural
ambiguity, because \texttt{10*1} padding is injective on field-sequence inputs. O3 adds only the
discipline that the domain string is globally unique to this statement.

\subsection*{O4: Zero-knowledge simulator (field-independent)}

\textbf{Implemented.} A hiding PCS (salted hiding-Merkle MMCS $+$ FRI on a vanishing-poly-randomized
codeword at \texttt{log\_blowup}\,$+1$) keeps the witness ($t$, $L$, siblings, direction bits, all
sponge intermediates) in the blinded trace; only $(\mathit{root}, \mathit{ctx}, N)$ is public. Two
blinding-randomness leaks that we found and fixed: (L1) the salts / blinding polynomials must come
from a CSPRNG, not a linear PRNG (a salted-but-linear generator is de-blindable when salts appear in
the proof); (L2) the seeds must be fresh, independent, $\geq 256$-bit per proof, drawn from an OS
CSPRNG, not fixed / guessable constants.

\textbf{The residual we seek review on.} A formal simulator $\mathcal{S}$ that, given only
$(\mathit{root}, \mathit{ctx}, N)$, outputs proofs computationally indistinguishable from honest
ones. The structure: $\mathcal{S}$ commits to a random masking polynomial of the same degree budget
and answers FRI / OOD queries via the standard ZK-STARK simulation \cite{BCRSVW19} $+$ hiding-PCS
zero-knowledge \cite{Hab22}; indistinguishability reduces to (i)~hiding of the salted Merkle
commitments and (ii)~the masking degree \textbf{strictly exceeding} the number of query $+$ OOD
openings, so every revealed evaluation is uniform and independent of the witness. The
\textbf{quantitative residual}: with the degree-$\alpha$ S-box forcing \texttt{log\_blowup}\,$\geq 3$ and a
concrete $\#\mathrm{queries} + \#\mathrm{OOD}$, the masking degree under the hiding PCS must be
confirmed $\geq \#\mathrm{openings} + 1$ for the production parameters, and the truncated-output
sponge openings ($\wout$ of $t$ cells) must leak nothing beyond the public digest under that budget.
O4 is field-independent; it transfers unchanged between Arm A and Arm B.

\subsection*{6.5 Proof-system (PCS / FRI) soundness: the challenge field}

O1--O4 concern the \emph{hash}; the proof system has its own soundness, bounded by the FRI / DEEP
challenge field and the query phase and binding on the Merkle commitment. The field-dependent terms
(out-of-domain / DEEP / proximity-gap) are \textbf{capped by the size of the challenge field},
regardless of query count, a subtlety that cost us a real error: Arm A originally drew its
challenges from its own $\approx 62$-bit value field $\GF{p^2}$, which caps DEEP soundness near $62$
bits, and a query-only estimate masked it. The fix is a larger challenge extension. Both arms now
compute to $128$-bit post-quantum soundness, \textbf{binding on a $256$-bit SHAKE256 Merkle
commitment} ($128$-bit collision, NIST Category 2):
\begin{itemize}[leftmargin=*]
\item \textbf{Arm A:} challenge field $\GF{p^6}$ (degree-3 over $\GF{p^2}$, $\approx 186$ bits),
  \texttt{log\_blowup}~3, $96$ queries, $20$ PoW bits $\to$ DEEP / Schwartz--Zippel $\approx 176$
  bits; binding term $128$ (SHAKE256).
\item \textbf{Arm B:} challenge field $\GF{q^5}$ (degree-5 over BabyBear, $\approx 155$ bits),
  \texttt{log\_blowup}~4, $96$ queries, $20$ PoW bits $\to$ DEEP / Schwartz--Zippel $\approx 147$
  bits; binding term $128$ (SHAKE256).
\end{itemize}
The post-quantum model is the deployed QROM one: Fiat--Shamir preserves a round-by-round-sound IOP up
to small factors, Grover halves only the grinding proof-of-work, and the SHAKE256 binding is the
$128$-bit term. This is a parameter characterization of the PCS layer; it assumes the AIR is sound
and complete (the structural audit and per-column mutation tests are evidence, not proof) and that
O1 holds.

\section{Scope and the questions we seek review on}

We do \textbf{not} claim O1--O4 are settled. The two residuals where independent expert review is
decisive:
\begin{enumerate}[leftmargin=*]
\item \textbf{(O1) The $\GF{p^2}$-\emph{state} parameterization (Arm A).} Running the \emph{entire
  permutation state} over $\GF{p^2}$ (rather than the base field, with the extension reserved for FRI
  challenges) is an unusual deployment choice. We have shown $8 + 60$ clears the published bounds by
  $\approx 2\times$ and that generic round constants close the subfield-descent path, but we seek
  confirmation of: (a)~the constant-genericity argument as the intended subfield defense; (b)~the
  absence of a residual invariant subspace under the base-field MDS $+$ generic-constant composition;
  (c)~comfort with the $\GF{p^2}$-state convention against current and foreseeable algebraic
  cryptanalysis. \emph{If a reviewer is uncomfortable, \textbf{Arm B, the recommended, implemented
  base-field instantiation, avoids this obligation entirely}, reducing O1 to conformance with
  published BabyBear Poseidon2 parameters.}
\item \textbf{(O4) The simulator's masking-degree budget.} We seek a formal simulator write-up and
  confirmation that the concrete masking degree $\geq \#\mathrm{queries} + \#\mathrm{OOD} + 1$ for the
  production parameters, and that the truncated-output sponge openings leak nothing under that budget.
\end{enumerate}

On \textbf{Arm A vs.\ Arm B}: we recommend Arm B on the measured merits (at equal $128$-bit security
its proofs are $\approx 14$--$18\times$ smaller and verify $\approx 16$--$17\times$ faster,
Section 5.4) \emph{and} on the cleaner O1 (a base-field parameter citation, not an extension-field
analysis). Arm A is the original $\GF{p^2}$ instantiation; we retain its full analysis because the
$\GF{p^2}$-state construction is of independent interest, but the off-envelope O1 above is why we do
not lead with it.

A reviewer wishing to reproduce the O1 evaluation needs only the parameters in Section 5 and the
public generators cited below; no proprietary material is involved.

\section{Related Work}

\begin{itemize}[leftmargin=*]
\item \textbf{Algebraic hashes.} Poseidon \cite{GKRRS21} and Poseidon2 \cite{GKS23} establish the
  round-count formulas and the base-field parameterizations we cite; the HADES design strategy and
  its algebraic cryptanalysis \cite{ABM23}, with the recent attack literature
  \cite{Bariant22,FreeLunch,CheapLunch,GKR25}, inform the O1 self-check.
\item \textbf{Sponge security.} The indifferentiability and truncated-output bounds are from
  Bertoni--Daemen--Peeters--Van Assche \cite{BDPV07,BDPV08}.
\item \textbf{Transparent STARKs.} FRI and STARKs \cite{BBHR18,BGKS20}; circle STARKs over
  Mersenne-31 \cite{HLP24}; zero-knowledge STARKs / hiding PCS \cite{BCRSVW19,Hab22}; the Plonky3
  toolkit \cite{Plonky3} provides the BabyBear Poseidon2 reference parameters (Arm B) and the
  two-adic FRI machinery both arms use.
\item \textbf{Post-quantum proof systems (the other family).} The principal alternative post-quantum
  ZK route is \emph{lattice-based}: LaBRADOR \cite{LaBRADOR} and the succinct-ZK toolkit of Biasioli
  et al.\ \cite{BBL26} (building on it) give sub-$100$\,KB proofs for general statements (the
  smallest among post-quantum schemes) at the cost of structured-lattice assumptions. We instead
  rest only on hash collision resistance (transparent, no algebraic structure), trading proof size
  for a more conservative assumption and a purpose-built membership relation.
\item \textbf{Anonymous membership / nullifiers.} Semaphore-style signaling \cite{Sem}, rate-limiting
  nullifiers \cite{RLN}, the PLUME deterministic-nullifier scheme \cite{PLUME}, and privacy-pool
  constructions establish the leaf / nullifier pattern; our
  contribution is a transparent, hash-only (post-quantum-oriented) instantiation over a small prime
  field, with the accompanying parameter analysis.
\end{itemize}

\section*{Artifact Availability}

All code, parameters, in-circuit audits, and the benchmarking harness behind the claims in this note
are open source, so that the construction (Sections 3--4), the parameter sets and security analysis
(Sections 5--6), and the measured comparison (Section 5.4) can be independently rebuilt and verified.

\begin{itemize}[leftmargin=*]
\item \textbf{Implementation.} \emph{Both arms} are implemented in the \texttt{lib-q-zkp} crate of the
  \texttt{libQ} library: \url{https://github.com/Enkom-Tech/libQ} (Apache-2.0), pinned at commit
  \texttt{34cb3f0} (tag \texttt{v0.0.8}). The shared in-circuit gadgets are \texttt{src/air/wide\_sponge.rs},
  \texttt{wide\_merkle\_path.rs}, and \texttt{wide\_hash.rs}; the membership API is
  \texttt{src/membership.rs} / \texttt{src/circuit.rs} / \texttt{src/api.rs}. \textbf{Arm A}
  (\(\GF{p^2}\) / Mersenne-31): membership AIR \texttt{src/air/unlinkable\_membership.rs}, gadget
  \texttt{src/air/poseidon\_gadget.rs} (\(\alpha=5\)), STARK config \texttt{src/stark.rs}, Merkle
  \texttt{src/merkle.rs}. \textbf{Arm B} (BabyBear / Poseidon2): membership AIR
  \texttt{src/air/unlinkable\_membership\_baby\_bear.rs}, gadget \texttt{src/air/poseidon2\_gadget.rs}
  (\(\alpha=7\)), STARK config \texttt{src/stark\_baby\_bear.rs}, Merkle
  \texttt{src/merkle\_baby\_bear.rs}. Both expose transparent and zero-knowledge prover/verifier paths.
\item \textbf{Reproduction package.} A companion repository bundles this paper's source together with
  a pinned snapshot (git submodule) of \texttt{libQ} and a \texttt{REPRODUCE.md} that maps each claim
  to the file, test, or benchmark that backs it:
  \url{https://github.com/En-Cipher/transparent-pq-anonymous-membership}.
\item \textbf{Security analysis (O1--O4).} The round-count / soundness derivations are recorded in
  \texttt{docs/membership-arm-a-soundness-params.md}, \texttt{membership-arm-b-soundness-params.md},
  \texttt{membership-arm-b-obligation-packet.md}, \texttt{membership-adr113-freeze-gate-review.md},
  and the adversarial self-check in \texttt{membership-arm-b-redteam.md}. The structural /
  under-constraint audit is exercised by \texttt{tests/air\_integration.rs},
  \texttt{tests/ip\_soundness\_tests.rs}, \texttt{tests/unlinkable\_membership\_tests.rs},
  \texttt{tests/zero\_knowledge\_tests.rs}, and \texttt{tests/security\_parameter\_tests.rs}.
\item \textbf{Measured cost (Section 5.4).} Produced by \texttt{benches/zkp\_benchmarks.rs} (harness
  \texttt{measure\_arm\_a} / Arm B), documented in \texttt{docs/membership-arm-b-measurement.md}.
  Environment: AMD Ryzen 9 5900X (12C/24T), Windows 10 IoT x86\_64, Rust 1.96 (edition 2024),
  \texttt{cargo test --release}, single-threaded (the \texttt{parallel}/rayon feature off), prove and
  verify reported as the median of 5, proof size via \texttt{postcard::to\_allocvec(proof).len()}.
  Absolute timings are hardware-dependent; the cross-arm \emph{ratios} are the reproducible result.
\end{itemize}

\section*{Acknowledgements}

\textbf{Use of generative AI tools.} The construction, the parameter choices, the security analysis,
and the open questions raised for review are the author's own work. Generative AI software (Anthropic
Claude, via Claude Code) was used as a drafting and editing aid for the exposition and to cross-check
the write-up against the cited literature; all technical content was directed, checked, and is
vouched for by the author. The author is solely responsible for the veracity and correctness of all
claims; generative AI is not, and cannot be, an author.

\begin{thebibliography}{Bariant+22}

\bibitem[GKRRS21]{GKRRS21} L.~Grassi, D.~Khovratovich, C.~Rechberger, A.~Roy, M.~Schofnegger.
  \emph{Poseidon: A New Hash Function for Zero-Knowledge Proof Systems.} USENIX Security 2021. IACR
  ePrint 2019/458.

\bibitem[GKS23]{GKS23} L.~Grassi, D.~Khovratovich, M.~Schofnegger. \emph{Poseidon2: A Faster Version
  of the Poseidon Hash Function.} AFRICACRYPT 2023. IACR ePrint 2023/323.

\bibitem[ABM23]{ABM23} T.~Ashur, T.~Buschman, M.~Mahzoun. \emph{Algebraic Cryptanalysis of the HADES
  Design Strategy: Application to Poseidon and Poseidon2.} ACISP 2024. IACR ePrint 2023/537.

\bibitem[BDPV07]{BDPV07} G.~Bertoni, J.~Daemen, M.~Peeters, G.~Van Assche. \emph{Sponge Functions.}
  ECRYPT Hash Workshop 2007.

\bibitem[BDPV08]{BDPV08} G.~Bertoni, J.~Daemen, M.~Peeters, G.~Van Assche. \emph{On the
  Indifferentiability of the Sponge Construction.} EUROCRYPT 2008, LNCS 4965, pp.~181--197.

\bibitem[BBHR18]{BBHR18} E.~Ben-Sasson, I.~Bentov, Y.~Horesh, M.~Riabzev. \emph{Scalable,
  Transparent, and Post-Quantum Secure Computational Integrity.} IACR ePrint 2018/046.

\bibitem[BGKS20]{BGKS20} E.~Ben-Sasson, L.~Goldberg, S.~Kopparty, S.~Saraf. \emph{DEEP-FRI: Sampling
  Outside the Box Improves Soundness.} ITCS 2020. IACR ePrint 2019/336.

\bibitem[HLP24]{HLP24} U.~Hab\"ock, D.~Levit, S.~Papini. \emph{Circle STARKs.} IACR ePrint 2024/278.

\bibitem[BCRSVW19]{BCRSVW19} E.~Ben-Sasson, A.~Chiesa, M.~Riabzev, N.~Spooner, M.~Virza, N.~Ward.
  \emph{Aurora: Transparent Succinct Arguments for R1CS.} EUROCRYPT 2019. IACR ePrint 2018/828.

\bibitem[Hab22]{Hab22} U.~Hab\"ock. \emph{A Summary on the FRI Low Degree Test.} IACR ePrint
  2022/1216.

\bibitem[Bariant+22]{Bariant22} A.~Bariant, C.~Bouvier, G.~Leurent, L.~Perrin. \emph{Algebraic
  Attacks against Some Arithmetization-Oriented Primitives.} IACR Trans.\ Symmetric Cryptol.\
  2022(3), pp.~73--101.

\bibitem[FreeLunch]{FreeLunch} A.~Bariant, A.~B\oe uf, A.~Lemoine, I.~Manterola Ayala, M.~\O ygarden,
  L.~Perrin, H.~Raddum. \emph{The Algebraic FreeLunch: Efficient Gr\"obner Basis Attacks Against
  Arithmetization-Oriented Primitives.} CRYPTO 2024. IACR ePrint 2024/347.

\bibitem[CheapLunch]{CheapLunch} \emph{The Algebraic CheapLunch: Extending FreeLunch Attacks on
  Arithmetization-Oriented Primitives Beyond CICO-1.} IACR ePrint 2025/2040.

\bibitem[GKR25]{GKR25} L.~Grassi, K.~Koschatko, C.~Rechberger. \emph{Poseidon and Neptune: Gr\"obner
  Basis Cryptanalysis Exploiting Subspace Trails.} IACR Trans.\ Symmetric Cryptol.\ 2025. IACR ePrint
  2025/954.

\bibitem[Plonky3]{Plonky3} The Plonky3 Authors. \emph{Plonky3: A toolkit for polynomial IOPs /
  STARKs.} Software, \url{https://github.com/Plonky3/Plonky3} (source of the BabyBear Poseidon2
  reference parameters).

\bibitem[LaBRADOR]{LaBRADOR} W.~Beullens, G.~Seiler. \emph{LaBRADOR: Compact Proofs for R1CS from
  Module-SIS.} CRYPTO 2023. IACR ePrint 2022/1341.

\bibitem[BBL+26]{BBL26} B.~Biasioli, M.~Bolboceanu, V.~Lyubashevsky, A.~Merino-Gallardo, M.~Osadnik,
  G.~Seiler, P.~Steuer. \emph{A Toolkit for Succinct Lattice-Based Zero-Knowledge Proofs.} IACR
  ePrint 2026/1289.

\bibitem[Sem]{Sem} Semaphore Protocol (Privacy \& Scaling Explorations). \emph{A zero-knowledge
  protocol for anonymous signaling / group membership.} Documentation:
  \url{https://docs.semaphore.pse.dev}; source: \url{https://github.com/semaphore-protocol/semaphore}.

\bibitem[RLN]{RLN} Rate-Limiting Nullifier (Privacy \& Scaling Explorations). \emph{A
  Merkle-membership $+$ nullifier spam-prevention gadget built on Semaphore (Shamir-secret-sharing key
  recovery on a rate-limit breach).} Documentation:
  \url{https://rate-limiting-nullifier.github.io/rln-docs/}; source:
  \url{https://github.com/Rate-Limiting-Nullifier}.

\bibitem[PLUME]{PLUME} A.~Gupta, K.~Gurkan. \emph{PLUME: An ECDSA Nullifier Scheme for Unique
  Pseudonymity within Zero-Knowledge Proofs.} IACR ePrint 2022/1255.

\end{thebibliography}

\end{document}
